\capitulo{4}{Técnicas y herramientas}

En este capítulo se describen las herramientas principales empleadas o consideradas durante la fase de pruebas del proyecto, centradas en el ámbito del procesamiento de imágenes y la visión por computador.  
Se presentan de forma general sus características, finalidad y fuentes oficiales de documentación.

\section{Herramientas fase de pruebas}

Durante la fase de pruebas del proyecto se han evaluado diversas herramientas para el procesamiento y análisis de imágenes.  
A continuación se describen las dos más relevantes: \textbf{MediaPipe} y \textbf{YOLO}, ambas ampliamente utilizadas en el campo de la visión artificial.

\subsection{MediaPipe}

MediaPipe es un marco de trabajo de código abierto desarrollado por Google para la creación de \textit{pipelines} o flujos de procesamiento multimedia en tiempo real \cite{lugaresi2019mediapipe, google_mediapipe_docs, google_mediapipe_github}.  

Permite construir sistemas modulares basados en grafos donde cada componente (denominado \textit{calculator}) realiza una tarea específica, como la captura de imágenes, el preprocesamiento, la inferencia de modelos de aprendizaje automático o el postprocesado de resultados.  
Su diseño facilita la reutilización de módulos y la integración de modelos de visión por computador en distintas plataformas (Android, iOS, escritorio o web).

Entre las soluciones más utilizadas que ofrece se incluyen:
\begin{itemize}
	\item Detección de rostro (\textit{Face Detection}) y malla facial 3D (\textit{Face Mesh}).
	\item Detección y seguimiento de manos (\textit{Hand Tracking}).
	\item Estimación de pose corporal completa (\textit{Pose}).
	\item Solución integral \textit{Holistic}, que combina rostro, manos y cuerpo.
\end{itemize}

Estas funcionalidades permiten aplicar MediaPipe a proyectos de reconocimiento facial, análisis de movimiento, interacción gestual, realidad aumentada y aplicaciones educativas o deportivas.

Toda la documentación oficial se encuentra disponible en la guía técnica de Google AI Edge \cite{google_mediapipe_docs}, y el código fuente está publicado en su repositorio oficial de GitHub \cite{google_mediapipe_github}.


\subsection{YOLO (You Only Look Once)}

YOLO es una familia de modelos de detección de objetos en tiempo real creada por la empresa Ultralytics.  
Su principio básico consiste en dividir una imagen en una cuadrícula y predecir simultáneamente las coordenadas y la clase de los objetos presentes en cada región \cite{ultralytics_yolov8, ultralytics_github}.

La versión más reciente, YOLOv8, incorpora arquitecturas modernas optimizadas para precisión y velocidad, con soporte para varias tareas de visión artificial:

\begin{itemize}
	\item Detección de objetos (\textit{object detection}).
	\item Segmentación de instancias (\textit{instance segmentation}).
	\item Clasificación de imágenes.
	\item Estimación de pose humana.
\end{itemize}

Se distribuye como paquete Python mediante la librería \ruta{ultralytics}, permitiendo la carga de modelos pre entrenados o el entrenamiento personalizado sobre conjuntos de datos propios.

La documentación y ejemplos de uso están disponibles en el sitio oficial de Ultralytics \cite{ultralytics_yolov8}, mientras que el código fuente se mantiene en su repositorio público de GitHub \cite{ultralytics_github}.



\subsection{Comparativa entre MediaPipe y YOLO}

Aunque ambos frameworks se utilizan en el ámbito de la visión por computador, su enfoque y finalidad son distintos y, en muchos casos, complementarios.  
La tabla siguiente resume sus principales diferencias:
\FloatBarrier
\tablaSmall{Comparación general entre MediaPipe y YOLO}{p{3cm}p{6cm}p{6cm}}{comparativaMPYOLO}{
	\textbf{Aspecto} & \textbf{MediaPipe} & \textbf{YOLO} \\
}{
	Enfoque principal & Procesamiento de características humanas (rostro, manos, pose). & Detección y localización de objetos genéricos en imágenes o vídeo. \\
	Tipo de modelo & Framework modular con soluciones preentrenadas de Google. & Red neuronal convolucional profunda (CNN) optimizada para detección. \\
	Ejecución & En tiempo real, incluso en dispositivos móviles o CPU. & En tiempo real, preferentemente con GPU para rendimiento óptimo. \\
	Configuración & Bajo nivel de configuración: soluciones listas para usar. & Permite entrenamiento y ajuste fino sobre datasets personalizados. \\
	Precisión en personas & Alta precisión en puntos clave del cuerpo humano. & Alta precisión en la localización de personas y objetos. \\
	Lenguaje y soporte & Implementación principal en Python y C++. & Implementación en Python con soporte mediante la librería Ultralytics. \\
	Licencia & Apache License 2.0 (Google). & AGPL-3.0 License (Ultralytics). \\
	\textbf{Uso ideal} & Análisis de movimiento, reconocimiento facial, interfaces gestuales. & Detección general de objetos, vigilancia, control de tráfico o inventario. \\
}

En resumen, \textbf{MediaPipe} resulta más adecuado para aplicaciones centradas en la estimación de características humanas y análisis del movimiento, mientras que \textbf{YOLO} se utiliza para la detección y clasificación de objetos en escenas más amplias.  
En conjunto, ambas herramientas pueden integrarse en sistemas híbridos donde se necesite combinar información semántica (qué objetos hay) y estructural (cómo se mueven o se relacionan las personas).


\subsection{Framework de servidor: evaluación y elección}

Para la capa de servidor del proyecto se necesita un API HTTP que reciba imágenes y parámetros, coordine los módulos de IA desarrollados en Python y asegure que podamos añadir después funciones de seguridad (autenticación, trazabilidad, CORS). Bajo estos objetivos, se han analizado tres frameworks.

\paragraph{FastAPI (Python).}  
\textbf{Ventajas:} permite programación asíncrona, genera automáticamente documentación interactiva, asegura tipos de datos con Pydantic.  
\textbf{Inconvenientes:} añade más estructura y definición de datos de la necesaria al inicio.  
\textbf{Motivo de descarte:} se buscaba avanzar rápido con una API sencilla y sin tanta formalidad estructural.

\paragraph{Django REST Framework (Python).}  
\textbf{Ventajas:} completo para proyectos grandes, con gestores, serializadores y autenticación ya integrada.  
\textbf{Inconvenientes:} demasiado pesado para un backend cuyo foco principal es orquestar inferencias de IA.  
\textbf{Motivo de descarte:} la complejidad no se ajusta al alcance de este TFG.

\paragraph{Express.js (Node.js).}  
\textbf{Ventajas:} minimalista, gran comunidad.  
\textbf{Inconvenientes:} implicaría usar otro lenguaje además de Python, lo que complica la integración con los modelos de IA.  
\textbf{Motivo de descarte:} se opta por mantener todo el servidor en Python para mayor coherencia.

\subsubsection{Framework elegido: Flask}

\paragraph{Descripción.}  
Flask es un microframework web en Python. Se caracteriza por ofrecer lo esencial —rutas, peticiones, respuestas— sin imponer una estructura compleja. Su diseño ligero permite añadir sólo lo que el proyecto necesite.

\paragraph{Ventajas para este proyecto.}  
\begin{itemize}
	\item Permite lanzar rápidamente un endpoint que reciba la imagen y devuelva un resultado, sin código innecesario.
	\item Al estar en Python facilita conectar directamente con los módulos de IA.
	\item Su simplicidad permite que, cuando lo necesitemos, añadamos seguridad, control de acceso o logging sin tener que rehacer la arquitectura.
	\item Los desarrollos iniciales son rápidos y el despliegue sencillo.
\end{itemize}

\paragraph{Justificación de elección.}  
Dado que el servidor funciona como una capa de orquestación —no como un sistema transaccional complejo— se requiere algo ligero, sencillo de mantener y que permita crecer cuando haga falta. Flask cumple claramente esos requisitos: arranque rápido, bajo coste inicial, buena integración con Python y evolución controlada.

\paragraph{Funcionamiento en el proyecto.}  
Se definirán rutas HTTP para subir imágenes o enviar parámetros, dichas rutas delegarán el trabajo a los módulos internos de IA y devolverán respuestas en formato JSON. La estructura del proyecto podrá separarse en rutas / servicios / módulos, lo que mantiene el código ordenado y escalable.
